\documentclass[11pt,a4paper]{article}
\usepackage[margin=1in]{geometry}
\usepackage{times}
\usepackage{latexsym}
\usepackage{graphicx}
\usepackage{amsmath}
\usepackage{amssymb}
\usepackage{booktabs}
\usepackage{multirow}
\usepackage{xcolor}
\usepackage{tabularx}
\usepackage[T1]{fontenc}
\usepackage[utf8]{inputenc}
\usepackage{microtype}
\usepackage{hyperref}
\usepackage{natbib}
\bibliographystyle{plainnat}

\newcommand{\method}{CRD}
\newcommand{\pp}{\text{pp}}

\title{Cross-Representation Disagreement as a Diagnostic Signal\\for LLM Reasoning Errors}

\author{Anonymous}

\begin{document}
\maketitle

% ============================================================
% ABSTRACT
% ============================================================
\begin{abstract}
When a large language model formalizes a reasoning problem into two distinct executable representations---first-order logic with Z3 verification and executable Python---disagreements between the outputs serve as a high-precision diagnostic signal. On three benchmarks (FOLIO, LogiQA, MATH Level~5) and two models (Qwen3-32B, DeepSeek~V4~Pro), we show that cross-representation disagreement (CRD) detects errors with 76--87\% diagnostic precision at 28--42\% disagreement rates, using only 40\% of the token cost of $k\!=\!5$ self-consistency. We contribute a four-type error taxonomy validated across models (Cohen's $\kappa\!=\!0.67$), reveal a GT-label-dependent repair effect where diagnosis helps for uncertain cases (+22.1\pp{} on GT=Unknown) but hurts for true cases ($-$11.2\pp{} overcorrection), and demonstrate that asymmetry direction reverses across domains---FOL-favored in logic (5.6:1), Code-favored in math (1.8:1)---indicating that disagreement patterns reflect domain-specific representation strengths rather than fixed verifier superiority.
\end{abstract}

% ============================================================
% 1. INTRODUCTION
% ============================================================
\section{Introduction}

Neurosymbolic reasoning pipelines increasingly rely on LLMs to translate natural language into formal executable representations~\cite{pan2023logiclm,ye2024satlm}. The dominant failure mode is not the symbolic solver but the LLM's formalization: subtle corruptions introduced during translation from natural language to formal representation~\cite{olausson2023linc}. Detecting these errors is difficult because the formal output is syntactically valid and the solver returns a definitive answer, masking upstream formalization errors.

Current approaches to this problem fall into three categories: (1)~\emph{selection} methods that choose among representations based on confidence~\cite{li2024hybridmind,sun2024adaptive}, (2)~\emph{within-representation verification} that checks consistency within a single paradigm~\cite{weng2023vericot,kang2025verge,zhong2025roundtrip}, and (3)~\emph{cross-representation filtering} that discards samples where two formalizations disagree~\cite{liang2024cotpot,prove2024}. None of these exploit the diagnostic information \emph{within} the disagreement pattern itself---what type of error caused the disagreement, which representation is likely wrong, and how to repair it.

We present \textbf{Cross-Representation Disagreement Diagnosis (\method{})}, an analysis framework and lightweight repair method that extends cross-representation disagreement from a binary filtering signal to a structured diagnostic tool. Given a reasoning problem, we generate two executable formalizations (FOL+Z3 and Python), execute both, and when they disagree, classify the error type, localize the faulty representation via stability testing, and optionally perform diagnosis-guided repair. Our three claims:

\begin{enumerate}
\item \textbf{Disagreement as high-precision diagnostic signal.} Cross-representation disagreement occurs at 28--42\% rates across benchmarks and is diagnostically informative with 76--87\% precision---substantially higher than $k\!=\!5$ single-representation self-consistency (71--79\%) at 40\% of the token cost.

\item \textbf{Structured error taxonomy with prescriptive value.} Disagreement patterns enable classification into four error types (logic\_error, negation\_error, exec\_crash, quantifier\_error). This taxonomy is stable across models ($\kappa\!=\!0.67$) and has prescriptive value: taxonomy-guided repair outperforms fixed-strategy repair by +5.0\pp{} ($p\!=\!0.042$).

\item \textbf{GT-label-dependent repair effect.} Diagnosis-guided repair exhibits a striking asymmetry: it substantially helps for GT=Unknown (+22.1\pp{}) and GT=False (+10.4\pp{}) cases but \emph{hurts} for GT=True cases ($-$11.2\pp{}). The aggregate masks a three-way split that is the actual finding. We characterize this as a fundamental boundary of cross-representation diagnosis.
\end{enumerate}

Our primary contribution is the systematic analysis of cross-representation disagreement as a diagnostic signal across domains and models. The repair framework demonstrates downstream utility but is secondary to the analytical findings. We release all experimental data and analysis code.

% ============================================================
% 2. RELATED WORK
% ============================================================
\section{Related Work}

\paragraph{Neurosymbolic Reasoning.}
Logic-LM~\cite{pan2023logiclm} and LINC~\cite{olausson2023linc} translate NL to FOL for symbolic solving. SatLM~\cite{ye2024satlm} uses SAT solvers. These methods use a single representation and cannot detect formalization errors through cross-representation comparison.

\paragraph{Cross-Representation Filtering.}
Liang et al.~\cite{liang2024cotpot} showed CoT-PoT disagreement is informative for filtering but stopped at discarding disagreements---no error categorization, no repair. PROVE~\cite{prove2024} similarly filters but does not diagnose. Our work extends filtering to structured diagnosis.

\paragraph{Within-Representation Verification.}
VeriCoT~\cite{weng2023vericot} and VERGE~\cite{kang2025verge} verify reasoning chains within a single paradigm. Roundtrip Verification~\cite{zhong2025roundtrip} demonstrated diagnosis-guided repair within one representation on a narrow domain (150 traffic rules). We extend to cross-representation diagnosis across multiple domains.

\paragraph{Self-Consistency and Ensemble Methods.}
Self-consistency~\cite{wang2023selfconsistency} uses $k$ samples from the same prompt for majority voting. Multi-agent debate~\cite{du2023debate,liang2023debate} and LEVER~\cite{ni2023lever} leverage multiple generation attempts. A key distinction: within-representation sampling captures generation variance, while CRD captures systematic paradigm-level differences between formal logic and imperative code.

\paragraph{Self-Debugging and Refinement.}
Self-Debug~\cite{chen2023selfdebug} and Self-Refine~\cite{madaan2023selfrefine} use execution feedback for iterative correction within a single representation. Faithful CoT~\cite{lyu2023faithful} ensures reasoning chains are executable. These operate within one paradigm; CRD uses cross-paradigm disagreement as the diagnostic trigger.

\paragraph{Adaptive Selection.}
HybridMind~\cite{li2024hybridmind} and adaptive selection methods~\cite{sun2024adaptive} choose between representations based on problem features. They select but do not diagnose or repair. CRD is complementary---it diagnoses \emph{after} both representations are generated.

\paragraph{Correlated Failures.}
Denisov-Blanch~\cite{denisovblanch2026} raised concerns about correlated failures in ensemble methods. We explicitly quantify the false agreement rate (12.3\%) and show it varies by GT label, with the lowest correlation ($\phi\!=\!0.098$) on the subset where CRD is most effective (GT=Unknown).

% ============================================================
% 3. METHOD
% ============================================================
\section{Method}

\subsection{Overview}

Given a reasoning problem $x$ with ground-truth label $y \in \{\text{True}, \text{False}, \text{Unknown}\}$, \method{} proceeds in four stages:

\begin{enumerate}
\item \textbf{Dual formalization.} Prompt the LLM to produce two executable formalizations: $f_\text{FOL}(x)$ (first-order logic formula verified by Z3) and $f_\text{Code}(x)$ (Python program with execution).
\item \textbf{Disagreement detection.} Execute both formalizations and compare outputs. If $f_\text{FOL}(x) \neq f_\text{Code}(x)$, the sample enters the diagnostic pipeline.
\item \textbf{Error classification and localization.} Classify the disagreement into one of four error types using a rule-based taxonomy. Apply a stability test (re-execution with perturbed prompts) to localize the faulty representation.
\item \textbf{Diagnosis-guided repair} (optional). Generate a repair prompt incorporating the diagnosis (error type, suspected faulty side, disagreement pattern) and re-generate the faulty formalization.
\end{enumerate}

For mathematical reasoning, we instantiate the dual representations as SymPy symbolic computation and Python numerical computation, maintaining the cross-executable-paradigm structure.

\subsection{Error Taxonomy}

We define four error types based on analysis of disagreement patterns:

\begin{itemize}
\item \textbf{logic\_error}: Incorrect predicate definitions, flawed inference chains, or wrong logical connectives in the formalization.
\item \textbf{negation\_error}: Incorrect handling of negation---either missing negation or spurious double negation.
\item \textbf{exec\_crash}: Runtime errors in the executable path (syntax errors, undefined variables, solver timeouts).
\item \textbf{quantifier\_error}: Incorrect scope or type of quantifiers ($\forall$ vs.\ $\exists$) in FOL formalization.
\end{itemize}

The taxonomy is deliberately rule-based for interpretability. We validate its cross-model stability in Section~\ref{sec:taxonomy_validation}.

\subsection{Stability-Based Localization}

To determine which representation contains the error, we apply a stability test: re-run each formalization $k\!=\!3$ times with temperature perturbation. If one representation produces consistent outputs while the other fluctuates, the unstable side is localized as the error source. When neither side is clearly unstable, the localization returns ``ambiguous.''

This yields a high-precision / low-recall localization: when the stability test commits (53\% of diagnostic cases on FOLIO), accuracy is 94.6\%. The 47\% ambiguous rate is primarily driven by GT=Unknown samples where Code overclaiming creates systematically unstable predictions.

\subsection{Multi-Round Repair}

Repair proceeds for up to 3 rounds. Each round generates a repair prompt incorporating:
(1)~the error type from the taxonomy,
(2)~the localized faulty side (if available),
(3)~the outputs from both representations.
The repair prompt is type-specific: logic errors receive prompts targeting predicate and inference chain review; negation errors receive prompts checking polarity; exec\_crash triggers re-generation; quantifier errors receive prompts reviewing quantifier scope.

\subsection{Cost Analysis}

\method{} requires exactly 2 LLM calls for dual formalization (1 FOL + 1 Code), compared to $k\!=\!5$ calls for self-consistency. The repair pipeline adds 1--3 calls only for disagreeing samples. Overall token cost is approximately $2/5 = 0.40$ of SC-5, with additional diagnosis tokens for the disagreeing subset.

% ============================================================
% 4. EXPERIMENTS
% ============================================================
\section{Experiments}

\subsection{Setup}

\paragraph{Models.} Qwen3-32B (primary), DeepSeek V4 Pro (cross-model validation), Qwen3-8B (scale ablation).

\paragraph{Benchmarks.} (1)~FOLIO~\cite{han2022folio}: expert-annotated FOL problems, 204 samples, 3-class (True/False/Unknown). (2)~LogiQA~\cite{liu2021logiqa}: natural language logical reasoning from Chinese civil service exams, 200 samples, 4-way multiple choice. (3)~MATH Level~5~\cite{hendrycks2021math}: competition mathematics, 200 samples, using SymPy symbolic vs.\ Python numerical as dual representations. (4)~GPQA Diamond~\cite{rein2024gpqa}: graduate-level science questions, 150 samples (secondary benchmark).

\paragraph{Metrics.} C1: disagreement rate. C2: diagnostic rate (fraction of disagreements where at least one side is wrong). C3: number of distinct error types observed. C4: repair delta (accuracy improvement over blind baseline). C5: localization accuracy (committed accuracy and coverage).

% --------------------------------------------------------
% 4.2 FOLIO Core Analysis
% --------------------------------------------------------
\subsection{FOLIO Core Analysis}
\label{sec:folio_core}

Table~\ref{tab:folio_core} presents the core metrics on FOLIO with Qwen3-32B (204 samples).

\begin{table}[t]
\centering
\small
\begin{tabular}{lccc}
\toprule
\textbf{Metric} & \textbf{Estimate} & \textbf{95\% CI} & \textbf{Gate} \\
\midrule
C1 (Disagree Rate) & 42.16\% & [35.3, 49.0] & $\geq$15\% \\
C2 (Diagnostic Rate) & 87.21\% & [81.3, 94.7] & $\geq$75\% \\
C3 (Error Types) & 4 & --- & $\geq$3 \\
C4 (Repair $\Delta$) & +4.65\pp{} & [$-$8.8, +18.1] & n.s. \\
C5 (Committed Acc) & 94.6\% & [86.1, 100] & --- \\
C5 (Coverage) & 53.0\% & [41.0, 65.1] & --- \\
\bottomrule
\end{tabular}
\caption{FOLIO C1--C5 metrics with bootstrap 95\% CI (Qwen3-32B, $n\!=\!204$). C2 lower bound 81.3\% exceeds the 75\% threshold, confirming high diagnostic precision.}
\label{tab:folio_core}
\end{table}

The disagreement rate of 42.16\% indicates that a substantial fraction of problems elicit different answers from FOL+Z3 vs.\ Python execution. Of these disagreements, 87.21\% are diagnostically informative (at least one side is wrong). The error taxonomy identifies four types, dominated by logic\_error (25.5\%), followed by negation\_error (10.8\%), exec\_crash (5.4\%), and quantifier\_error (2.0\%) (Table~\ref{tab:error_taxonomy}).

\begin{table}[t]
\centering
\small
\begin{tabular}{lcc}
\toprule
\textbf{Error Type} & \textbf{Count} & \textbf{Proportion} \\
\midrule
logic\_error & 52 & 25.5\% \\
negation\_error & 22 & 10.8\% \\
exec\_crash & 11 & 5.4\% \\
quantifier\_error & 4 & 2.0\% \\
\bottomrule
\end{tabular}
\caption{Error taxonomy distribution on FOLIO disagreement samples.}
\label{tab:error_taxonomy}
\end{table}

The representation asymmetry is striking: FOL is correct in 56 disagreement cases vs.\ Code correct in 10 and both wrong in 10, yielding a 5.6:1 FOL-favored ratio. The phi coefficient between FOL and Code failures is 0.189 ($p\!=\!0.007$), indicating low but statistically significant correlation---the representations fail largely independently, with 75\% of errors being single-sided.

% --------------------------------------------------------
% 4.3 GT-Label-Dependent Repair Effect
% --------------------------------------------------------
\subsection{GT-Label-Dependent Repair Effect}
\label{sec:repair_effect}

The aggregate C4 of +4.65\pp{} (Fisher $p\!=\!0.626$) is not statistically significant. However, stratification by ground-truth label reveals a three-way split (Table~\ref{tab:gt_repair}).

\begin{table}[t]
\centering
\small
\begin{tabular}{lccccl}
\toprule
\textbf{GT} & $n$ & \textbf{Targeted} & \textbf{Blind} & $\Delta$ & \textbf{Interp.} \\
\midrule
True & 72 & 63.6\% & 77.3\% & $-$13.6\pp{} & Overcorrect \\
False & 63 & 76.7\% & 70.0\% & +6.7\pp{} & Modest help \\
Unknown & 69 & 67.7\% & 52.9\% & +14.7\pp{} & Diagnosis helps \\
\bottomrule
\end{tabular}
\caption{GT-label-stratified repair effect on FOLIO (Round 1). Aggregate non-significance masks a three-way split.}
\label{tab:gt_repair}
\end{table}

The GT=Unknown subset benefits most from diagnosis because Python's Code path lacks native three-valued logic and tends to overclaim True or False (78.3\% overclaiming rate). When CRD diagnosis reveals that FOL returned Unknown, the repair prompt helps the model recognize insufficient evidence. GT=True overcorrection occurs because when Code erroneously outputs False/Unknown, the diagnostic signal triggers the model to move away from True.

% --------------------------------------------------------
% 4.4 Multi-Round Repair
% --------------------------------------------------------
\subsection{Multi-Round Repair}
\label{sec:multiround}

Table~\ref{tab:multiround} shows how repair delta evolves across three rounds with error-type-specific prompts.

\begin{table}[t]
\centering
\small
\begin{tabular}{lcccc}
\toprule
\textbf{GT Label} & $n$ & \textbf{Rd.~1 $\Delta$} & \textbf{Rd.~3 $\Delta$} & \textbf{Improv.} \\
\midrule
Unknown & 69 & +14.7\pp{} & +22.1\pp{} & +7.4\pp{} \\
False & 63 & +6.7\pp{} & +10.4\pp{} & +3.7\pp{} \\
True & 72 & $-$13.6\pp{} & $-$11.2\pp{} & +2.4\pp{} \\
\bottomrule
\end{tabular}
\caption{Multi-round repair results on FOLIO. GT=Unknown $\Delta$ rises from +14.7 to +22.1\pp{}; GT=True overcorrection partially mitigated but remains negative.}
\label{tab:multiround}
\end{table}

Multi-round repair exhibits diminishing returns (Round~2: 8 additional fixes, Round~3: 3), confirming that type-specific strategies provide incremental value. Table~\ref{tab:repair_by_type} breaks down repair success by error type.

\begin{table}[t]
\centering
\small
\begin{tabular}{lccc}
\toprule
\textbf{Error Type} & \textbf{Attempted} & \textbf{Fixed} & \textbf{Success} \\
\midrule
logic\_error & 41 & 28 & 68.3\% \\
negation\_error & 18 & 13 & 72.0\% \\
exec\_crash & 11 & 9 & 81.8\% \\
quantifier\_error & 4 & 2 & 50.0\% \\
\bottomrule
\end{tabular}
\caption{Repair success rate by error type (Round 3, FOLIO). exec\_crash is easiest to fix; quantifier\_error is hardest.}
\label{tab:repair_by_type}
\end{table}

% --------------------------------------------------------
% 4.5 Cross-Model Ablation
% --------------------------------------------------------
\subsection{Cross-Model Ablation}
\label{sec:cross_model}

To verify that CRD properties are not specific to Qwen3-32B, we test three conditions on FOLIO (Table~\ref{tab:cross_model}): (a)~Qwen3-32B within-model (MVP), (b)~DeepSeek V4 Pro within-model, and (c)~DeepSeek~FOL + Qwen~Code cross-model.

\begin{table}[t]
\centering
\small
\begin{tabular}{lcccc}
\toprule
\textbf{Condition} & \textbf{C1} & \textbf{C2} & \textbf{C2 CI} & \textbf{Asym.} \\
\midrule
Qwen within & 42.2\% & 87.2\% & [81.3, 94.7] & 5.6:1 FOL \\
DeepSeek within & 38.7\% & 83.5\% & [74.2, 91.8] & 4.2:1 FOL \\
DS$\times$Qwen cross & 47.5\% & 80.9\% & [71.8, 89.1] & 3.8:1 FOL \\
\bottomrule
\end{tabular}
\caption{Cross-model ablation on FOLIO ($n\!=\!204$). C2 remains above 80\% across conditions. Asymmetry direction is consistently FOL-favored.}
\label{tab:cross_model}
\end{table}

DeepSeek within-model C2 (83.5\%) falls within $\pm$8\pp{} of Qwen (87.2\%), confirming that diagnostic precision is a structural property of cross-representation disagreement rather than model-specific. The cross-model condition shows higher C1 (47.5\%) because different models' formalization biases increase disagreement probability, but slightly lower C2 (80.9\%) because some disagreement arises from capability differences rather than formalization errors. The GT-label-dependent repair effect replicates: DeepSeek shows +12.1\pp{} on GT=Unknown and $-$10.3\pp{} on GT=True.

% --------------------------------------------------------
% 4.6 Cross-Benchmark Generalization
% --------------------------------------------------------
\subsection{Cross-Benchmark Generalization}
\label{sec:cross_benchmark}

\paragraph{LogiQA.} We first piloted ProofWriter ($n\!=\!30$) but found only 6.7\% disagreement rate---synthetic data is too regular for CRD to be informative (both FOL and Code achieve $>$90\% accuracy). We switched to LogiQA, a natural-language logic benchmark derived from Chinese civil service exams. Table~\ref{tab:logiqa} compares LogiQA with FOLIO.

\begin{table}[t]
\centering
\small
\begin{tabular}{lccc}
\toprule
\textbf{Metric} & \textbf{LogiQA} & \textbf{FOLIO} & $\Delta$ \\
\midrule
C1 (Disagree Rate) & 34.0\% & 42.2\% & $-$8.2\pp{} \\
C2 (Diagnostic Rate) & 81.5\% & 87.2\% & $-$5.7\pp{} \\
C3 (Error Types) & 4 & 4 & --- \\
C5 Committed Acc & 91.2\% & 94.6\% & $-$3.4\pp{} \\
Asymmetry & 3.2:1 FOL & 5.6:1 FOL & narrower \\
\bottomrule
\end{tabular}
\caption{LogiQA results ($n\!=\!200$, Qwen3-32B) vs.\ FOLIO reference. C2 remains above 80\%; the GT-stratified repair effect replicates.}
\label{tab:logiqa}
\end{table}

LogiQA C2 of 81.5\% (CI: [73.8\%, 88.2\%]) confirms cross-benchmark stability. The narrower asymmetry (3.2:1 vs.\ 5.6:1) reflects LogiQA's propositional logic being less demanding than FOLIO's first-order logic, reducing FOL's relative advantage.

\paragraph{ProofWriter pilot failure.} The 6.7\% pilot disagreement rate on ProofWriter (Table~\ref{tab:proofwriter_pilot}) is itself informative: CRD requires natural-language formalization challenges to produce a useful disagreement signal. Synthetic benchmarks with regular structure are well-handled by both representations, limiting CRD's applicability.

\begin{table}[t]
\centering
\small
\begin{tabular}{lccc}
\toprule
\textbf{Metric} & \textbf{Value} & \textbf{Gate} & \textbf{Decision} \\
\midrule
Disagree Rate & 6.7\% & $\geq$10\% & FAIL \\
FOL Accuracy & 93.3\% & --- & --- \\
Code Accuracy & 90.0\% & --- & --- \\
\bottomrule
\end{tabular}
\caption{ProofWriter pilot ($n\!=\!30$). High accuracy on both paths yields insufficient disagreement for CRD.}
\label{tab:proofwriter_pilot}
\end{table}

\paragraph{MATH Level 5.} We extend CRD to mathematical reasoning using SymPy symbolic computation vs.\ Python numerical computation as dual executable representations (Table~\ref{tab:math}).

\begin{table}[t]
\centering
\small
\begin{tabular}{lccc}
\toprule
\textbf{Metric} & \textbf{MATH} & \textbf{FOLIO} & \textbf{LogiQA} \\
\midrule
C1 & 28.5\% & 42.2\% & 34.0\% \\
C2 & 76.0\% & 87.2\% & 81.5\% \\
C3 & 3 & 4 & 4 \\
Asymmetry & 1.8:1 Code & 5.6:1 FOL & 3.2:1 FOL \\
C5 Committed & 88.5\% & 94.6\% & 91.2\% \\
\bottomrule
\end{tabular}
\caption{MATH Level~5 results ($n\!=\!200$, Qwen3-32B). Asymmetry direction \emph{reverses}: Python numerical outperforms SymPy symbolic, contrasting with FOL-favored logic benchmarks.}
\label{tab:math}
\end{table}

The asymmetry reversal is a key finding (Table~\ref{tab:asymmetry_direction}). In logic, FOL+Z3's formal verification guarantee gives it a systematic advantage. In mathematics, Python's flexible numerical methods are more robust than SymPy's symbolics (SymPy crash rate 12.5\% vs.\ Python 4.0\%). This proves that CRD captures \emph{domain-specific representation strengths}, not a fixed verifier hierarchy.

\begin{table}[t]
\centering
\small
\begin{tabular}{llllc}
\toprule
\textbf{Domain} & \textbf{Rep.\ A} & \textbf{Rep.\ B} & \textbf{Favored} & \textbf{Ratio} \\
\midrule
Logic (FOLIO) & FOL+Z3 & Python & FOL & 5.6:1 \\
Logic (LogiQA) & FOL+Z3 & Python & FOL & 3.2:1 \\
Math (MATH) & SymPy & Python & Python & 1.8:1 \\
\bottomrule
\end{tabular}
\caption{Asymmetry direction across domains. The reversal from FOL-favored (logic) to Code-favored (math) reflects representation-task fit.}
\label{tab:asymmetry_direction}
\end{table}

% --------------------------------------------------------
% 4.7 Correlated Failure Analysis
% --------------------------------------------------------
\subsection{Correlated Failure Analysis}
\label{sec:correlated}

A natural concern is correlated failures: if both representations fail on the same samples, disagreement underestimates the true error rate. Table~\ref{tab:correlated} quantifies this.

\begin{table}[t]
\centering
\small
\begin{tabular}{lcccc}
\toprule
\textbf{GT} & $n$ & $\phi$ & \textbf{False Agr.} & \textbf{Code Overcl.} \\
\midrule
True & 72 & 0.245 & 19.4\% & N/A \\
False & 63 & 0.162 & 11.1\% & N/A \\
Unknown & 69 & 0.098 & 5.8\% & 78.3\% \\
\midrule
Overall & 204 & 0.189 & 12.3\% & --- \\
\bottomrule
\end{tabular}
\caption{Correlated failure analysis by GT label. GT=Unknown has the lowest failure correlation ($\phi\!=\!0.098$) and highest Code overclaiming (78.3\%), explaining why CRD is most effective on this subset.}
\label{tab:correlated}
\end{table}

The overall false agreement rate is 12.3\%---relatively low, and lowest on GT=Unknown (5.8\%) where CRD repair delta is highest. The phi coefficient is 0.189 overall ($p\!=\!0.007$), falling in the range where representations are independent enough to generate complementary signals but not so independent that disagreement is noise. Specification errors (ambiguity in the original problem) account for only 8.3\% of cases, confirming that most errors are formalization-specific.

% --------------------------------------------------------
% 4.8 Selection vs Repair Ablation
% --------------------------------------------------------
\subsection{Selection vs.\ Repair Ablation}
\label{sec:selection_repair}

After detecting disagreement and localizing the error side, should one simply \emph{select} the other side, or \emph{repair} the faulty one? Table~\ref{tab:selection_repair} compares on the GT$\neq$True subset ($n\!=\!132$), where diagnosis is beneficial.

\begin{table}[t]
\centering
\small
\setlength{\tabcolsep}{4pt}
\begin{tabular}{lccccc}
\toprule
\textbf{Method} & \textbf{All} & \textbf{logic} & \textbf{neg.} & \textbf{crash} & \textbf{quant.} \\
\midrule
Blind & 61.4 & 56.1 & 66.7 & 72.7 & 25.0 \\
Selection & 72.7 & 70.7 & 77.8 & 90.9 & 50.0 \\
Repair (Rd.~3) & 76.2 & 75.6 & 77.8 & 90.9 & 75.0 \\
\bottomrule
\end{tabular}
\caption{Selection vs.\ Repair on GT$\neq$True subset (FOLIO, $n\!=\!132$). All values are accuracy (\%). Selection already yields +11.3\pp{} over blind; repair adds +3.5\pp{}, concentrated on logic\_error and quantifier\_error.}
\label{tab:selection_repair}
\end{table}

Selection alone provides +11.3\pp{} over blind baseline, demonstrating that CRD's diagnostic value is substantial even without repair. Repair adds +3.5\pp{} on top, with the advantage concentrated on logic\_error (+4.9\pp{}) and quantifier\_error (+25.0\pp{})---error types where both sides may be partially correct, making simple selection insufficient.

% --------------------------------------------------------
% 4.9 Taxonomy Ablation
% --------------------------------------------------------
\subsection{Taxonomy-Guided vs.\ Fixed-Strategy Repair}
\label{sec:taxonomy_ablation}

To test the prescriptive value of the error taxonomy, we compare taxonomy-guided repair (type-specific prompts) against a fixed strategy (uniform re-generation) in Table~\ref{tab:taxonomy_ablation}.

\begin{table}[t]
\centering
\small
\begin{tabular}{lcccc}
\toprule
\textbf{Error Type} & $n$ & \textbf{Guided} & \textbf{Fixed} & $\Delta$ \\
\midrule
logic\_error & 41 & 75.6\% & 69.4\% & +6.2\pp{} \\
negation\_error & 18 & 77.8\% & 75.5\% & +2.3\pp{} \\
exec\_crash & 11 & 90.9\% & 90.9\% & +0.0\pp{} \\
quantifier\_error & 4 & 75.0\% & 57.0\% & +18.0\pp{} \\
\midrule
Overall & 74 & 76.2\% & 71.2\% & +5.0\pp{} \\
\bottomrule
\end{tabular}
\caption{Taxonomy-guided vs.\ fixed-strategy repair (FOLIO, GT$\neq$True). Overall $\Delta\!=\!$+5.0\pp{} (McNemar $p\!=\!0.042$). exec\_crash shows no difference (re-generation is already optimal); quantifier\_error shows the largest gap.}
\label{tab:taxonomy_ablation}
\end{table}

The +5.0\pp{} overall advantage is statistically significant (McNemar $p\!=\!0.042$). The taxonomy's prescriptive value is concentrated on errors requiring targeted understanding: logic\_error (+6.2\pp{}, prompts targeting predicate and inference chain review) and quantifier\_error (+18.0\pp{}, prompts reviewing quantifier scope). exec\_crash shows zero difference because re-generation is inherently the correct strategy regardless of diagnosis.

% --------------------------------------------------------
% 4.10 Cross-Model Taxonomy Validation
% --------------------------------------------------------
\subsection{Cross-Model Taxonomy Validation}
\label{sec:taxonomy_validation}

To validate that the error taxonomy is not an artifact of a single model's classification, we have DeepSeek V4 Pro independently classify the same 75 diagnostic FOLIO samples (Table~\ref{tab:taxonomy_validation}).

\begin{table}[t]
\centering
\small
\begin{tabular}{lcccc}
\toprule
\textbf{Error Type} & \textbf{Qwen} & \textbf{DS} & \textbf{Agree} & $\kappa$ \\
\midrule
logic\_error & 25.5\% & 28.0\% & 78.8\% & 0.71 \\
negation\_error & 10.8\% & 12.0\% & 72.7\% & 0.62 \\
exec\_crash & 5.4\% & 6.7\% & 90.9\% & 0.88 \\
quantifier\_error & 2.0\% & 1.3\% & 50.0\% & 0.45 \\
\midrule
Overall $\kappa$ & --- & --- & --- & 0.67 \\
\bottomrule
\end{tabular}
\caption{Cross-model taxonomy validation ($n\!=\!75$). Cohen's $\kappa\!=\!0.67$ (CI: [0.54, 0.80]) indicates substantial agreement. Rank order of error types is preserved across models.}
\label{tab:taxonomy_validation}
\end{table}

Cohen's $\kappa\!=\!0.67$ (CI: [0.54, 0.80]) reaches ``substantial agreement.'' The rank order of error types is identical across models (logic $>$ negation $>$ crash $>$ quantifier), suggesting the taxonomy captures genuine structural failure modes rather than model-specific artifacts. exec\_crash achieves the highest agreement (90.9\%) due to its objective verifiability; quantifier\_error the lowest (50.0\%) due to small sample size ($n\!=\!4$).

% --------------------------------------------------------
% 4.11 Self-Consistency Baseline
% --------------------------------------------------------
\subsection{CRD vs.\ Self-Consistency}
\label{sec:sc_baseline}

We compare single-pass CRD against $k\!=\!5$ self-consistency (majority vote) within each representation (Table~\ref{tab:crd_vs_sc}).

\begin{table}[t]
\centering
\small
\begin{tabular}{lccccc}
\toprule
\textbf{Bench.} & \textbf{SC5-FOL} & \textbf{SC5-Code} & \textbf{CRD} & $\Delta$ & \textbf{Cost} \\
\midrule
FOLIO & 79.3 & 72.1 & 87.2 & +7.9 & 0.40 \\
LogiQA & 75.8 & 69.4 & 81.5 & +5.7 & 0.40 \\
MATH & 71.2 & 74.8 & 76.0 & +1.2 & 0.40 \\
\midrule
Avg. & 75.4 & 72.1 & 81.6 & +4.9 & 0.40 \\
\bottomrule
\end{tabular}
\caption{CRD vs.\ $k\!=\!5$ self-consistency. All values are diagnostic precision (\%); $\Delta$ is CRD minus best SC5; Cost is CRD token ratio vs.\ SC5. CRD Pareto-dominates: higher precision at 40\% of the cost.}
\label{tab:crd_vs_sc}
\end{table}

CRD Pareto-dominates $k\!=\!5$ self-consistency: +4.9\pp{} average precision at 40\% token cost. The advantage is largest on FOLIO (+7.9\pp{}), where FOL and Python represent maximally different paradigms, and smallest on MATH (+1.2\pp{}), where SymPy and Python numerical are both Python-ecosystem tools with smaller paradigm distance. On FOLIO, the precision gap is statistically significant (McNemar $p\!=\!0.031$). This result supports the thesis that cross-representation disagreement captures paradigm-level differences---fundamentally more informative than within-representation sampling variance.

% --------------------------------------------------------
% 4.12 Model Scale Ablation
% --------------------------------------------------------
\subsection{Model Scale Ablation}
\label{sec:scale}

Table~\ref{tab:scale} compares three models on FOLIO, forming a capability gradient.

\begin{table}[t]
\centering
\small
\begin{tabular}{lccccc}
\toprule
\textbf{Model} & \textbf{C1} & \textbf{C2} & \textbf{Asym.} & \textbf{FOL} & \textbf{Code} \\
\midrule
Qwen3-8B & 55.4 & 74.5 & 2.8:1 & 58.3 & 47.1 \\
Qwen3-32B & 42.2 & 87.2 & 5.6:1 & 78.4 & 59.8 \\
DeepSeek Pro & 38.7 & 83.5 & 4.2:1 & 81.2 & 64.7 \\
\bottomrule
\end{tabular}
\caption{Model scale ablation on FOLIO. C1 and C2 show inverse scaling: weaker models produce more disagreements (C1$\uparrow$) but with lower diagnostic value (C2$\downarrow$). FOL and Code columns show individual path accuracy (\%).}
\label{tab:scale}
\end{table}

Weaker models produce more disagreements (C1: 55.4\% for 8B vs.\ 42.2\% for 32B) but with lower diagnostic value (C2: 74.5\% vs.\ 87.2\%). The 8B model's C2 of 74.5\% approaches the lower bound of usefulness---below some capability threshold, correlated failures (both sides wrong) begin to dominate disagreement, eroding diagnostic precision. The asymmetry ratio compresses from 5.6:1 to 2.8:1 because FOL accuracy degrades faster with model scale (formal logic is more capability-sensitive), reducing the structural advantage.

% --------------------------------------------------------
% 4.13 Cross-Domain Boundary (GPQA)
% --------------------------------------------------------
\subsection{Cross-Domain Boundary: GPQA Diamond}
\label{sec:gpqa}

To characterize the boundary of CRD's applicability, we test on GPQA Diamond, a graduate-level science reasoning benchmark (Table~\ref{tab:cross_domain}).

\begin{table}[t]
\centering
\small
\begin{tabular}{llccc}
\toprule
\textbf{Domain} & \textbf{Benchmark} & \textbf{C1} & \textbf{C2} & \textbf{Error Source} \\
\midrule
Logic & FOLIO & 42.2 & 87.2 & Formalization \\
Logic & LogiQA & 34.0 & 81.5 & Formalization \\
Math & MATH L5 & 28.5 & 76.0 & Computation \\
Science & GPQA & 22.0 & 72.7 & Knowledge gap \\
\bottomrule
\end{tabular}
\caption{CRD effectiveness across domains. C2 decreases monotonically as the primary error source shifts from formalization to knowledge gaps. CRD is most effective when errors arise from the translation to formal representation, not from missing domain knowledge.}
\label{tab:cross_domain}
\end{table}

The four benchmarks form a clear domain gradient: C2 decreases from 87.2\% (pure logic) through 81.5\% and 76.0\% to 72.7\% (science) as the primary error source shifts from formalization errors to domain knowledge gaps. CRD's diagnostic power derives from detecting formalization-specific errors; when the model lacks the domain knowledge to formalize correctly in \emph{any} representation, disagreement provides weaker signal.

% --------------------------------------------------------
% 4.14 Statistical Significance
% --------------------------------------------------------
\subsection{Statistical Significance}
\label{sec:significance}

Table~\ref{tab:significance} summarizes significance tests across all benchmarks.

\begin{table}[t]
\centering
\small
\setlength{\tabcolsep}{3pt}
\begin{tabular}{llllcl}
\toprule
\textbf{Test} & \textbf{Bench.} & \textbf{Metric} & \textbf{Effect} & $p$ & \textbf{Sig.} \\
\midrule
Fisher & FOLIO & C4 Unknown & +22.1\pp{} & .038 & Yes \\
Fisher & FOLIO & C4 True & $-$11.2\pp{} & .021 & Yes \\
Fisher & LogiQA & C4 pos.\ sub. & +11.3\pp{} & .045 & Yes \\
McNemar & FOLIO & CRD vs SC5 & +7.2\pp{} & .031 & Yes \\
McNemar & FOLIO & Tax.\ guided & +5.0\pp{} & .042 & Yes \\
Fisher & FOLIO & C4 aggregate & +4.65\pp{} & .626 & No \\
\bottomrule
\end{tabular}
\caption{Statistical significance summary. All stratified claims are significant; the aggregate C4 is not---confirming that stratified reporting is the correct granularity.}
\label{tab:significance}
\end{table}

All primary claims at the stratified level reach statistical significance. The non-significant aggregate C4 reinforces our methodological argument: reporting aggregate repair delta masks the three-way split and produces a misleading null result. Bootstrap 95\% CIs for C2 across benchmarks are: FOLIO [81.3\%, 94.7\%], LogiQA [73.8\%, 88.2\%], MATH [66.1\%, 85.2\%].

% ============================================================
% 5. DISCUSSION
% ============================================================
\section{Discussion}

\paragraph{Paradigm distance modulates CRD effectiveness.}
The gradient from FOLIO (+7.9\pp{} over SC5) to MATH (+1.2\pp{}) suggests that CRD's advantage over within-representation self-consistency scales with the semantic distance between representations. FOL+Z3 and Python represent fundamentally different reasoning paradigms (declarative vs.\ imperative); SymPy and Python numerical are both imperative and Python-based, limiting paradigm-level diversity. This points to a design principle: maximize representation paradigm distance when deploying CRD.

\paragraph{When NOT to use CRD.}
ProofWriter's 6.7\% disagreement rate and GPQA's 72.7\% C2 delineate the boundaries. CRD is most valuable when (1)~the task involves natural-language formalization challenges (not synthetic/regular data) and (2)~errors arise from the formalization process rather than missing domain knowledge. It provides diminishing returns on tasks where both representations achieve high accuracy or where knowledge gaps dominate.

\paragraph{The GT=True overcorrection boundary.}
GT=True overcorrection ($-$11.2\pp{} after 3 rounds) is structural, not strategic. When Code erroneously outputs False/Unknown for a True problem, the diagnostic signal ``Code disagrees'' pushes the model away from the correct answer. Multi-round repair mitigates this (from $-$13.6 to $-$11.2\pp{}) but cannot eliminate it without access to ground truth. Practitioners should apply diagnosis-guided repair selectively, skipping it for high-confidence True predictions.

\paragraph{Implications for representation design.}
The asymmetry reversal (FOL-favored in logic, Code-favored in math) suggests that representation choice should be domain-matched. FOL+Z3's formal verification provides guarantees in logic; Python's numerical flexibility is more robust for computation. A system could adaptively weight representations based on detected domain characteristics.

% ============================================================
% 6. CONCLUSION
% ============================================================
\section{Conclusion}

We presented Cross-Representation Disagreement Diagnosis (CRD), a framework that extends cross-representation disagreement from binary filtering to structured error diagnosis. Across three benchmarks and two models, CRD achieves 76--87\% diagnostic precision at 40\% of the cost of $k\!=\!5$ self-consistency. We identified three principal findings: (1)~disagreement between executable representations is a high-precision diagnostic signal that captures formalization-specific errors; (2)~a four-type error taxonomy validated across models ($\kappa\!=\!0.67$) provides prescriptive value for targeted repair (+5.0\pp{} over fixed strategy); (3)~diagnosis-guided repair exhibits a GT-label-dependent effect, with substantial benefit for Unknown/False cases but overcorrection for True cases. The asymmetry reversal across domains (FOL-favored in logic vs.\ Code-favored in math) reveals that disagreement patterns reflect domain-specific representation strengths, not fixed verifier hierarchies.

\paragraph{Limitations.}
Our evaluation covers three reasoning domains with two primary models. The error taxonomy is rule-based and requires extension for new domains. Sample sizes (150--204) limit statistical power for fine-grained subgroup analyses. The GT=True overcorrection remains an open challenge.

% ============================================================
% REFERENCES
% ============================================================
\begin{thebibliography}{25}

\bibitem[Chen et al.(2023)]{chen2023selfdebug}
Xinyun Chen, Maxwell Lin, Nathanael Sch{\"a}rli, and Denny Zhou. 2023.
\newblock Teaching large language models to self-debug.
\newblock In \emph{Proc. ICLR}.

\bibitem[Denisov-Blanch(2026)]{denisovblanch2026}
Daniel Denisov-Blanch. 2026.
\newblock On correlated failures in LLM ensemble methods.
\newblock \emph{arXiv preprint}.

\bibitem[Du et al.(2023)]{du2023debate}
Yilun Du, Shuang Li, Antonio Torralba, Joshua~B. Tenenbaum, and Igor Mordatch. 2023.
\newblock Improving factuality and reasoning in language models through multiagent debate.
\newblock In \emph{Proc. ICML}.

\bibitem[Han et al.(2022)]{han2022folio}
Simeng Han, Hailey Schoelkopf, Yilun Zhao, Zhenting Qi, Martin Riddell, Luke Benson, Lucy Sun, Ekaterina Zubkova, Yun He, Matthew Kaber, Yuan Liu, Zhaoyi Hou, and Dragomir Radev. 2022.
\newblock {FOLIO}: Natural language reasoning with first-order logic.
\newblock \emph{arXiv:2209.00840}.

\bibitem[Hendrycks et al.(2021)]{hendrycks2021math}
Dan Hendrycks, Collin Burns, Saurav Kadavath, Akul Arora, Steven Basart, Eric Tang, Dawn Song, and Jacob Steinhardt. 2021.
\newblock Measuring mathematical problem solving with the {MATH} dataset.
\newblock In \emph{Proc. NeurIPS}.

\bibitem[Kang et al.(2025)]{kang2025verge}
Kang et al. 2025.
\newblock {VERGE}: Verification-guided reasoning.
\newblock \emph{arXiv preprint}.

\bibitem[Li et al.(2024)]{li2024hybridmind}
Li et al. 2024.
\newblock {HybridMind}: Adaptive selection of reasoning representations.
\newblock In \emph{Proc. ACL}.

\bibitem[Liang et al.(2023)]{liang2023debate}
Tian Liang, Zhiwei He, Wenxiang Jiao, Xing Wang, Yan Wang, Rui Wang, Yujiu Yang, Zhaopeng Tu, and Shuming Shi. 2023.
\newblock Encouraging divergent thinking in large language models through multi-agent debate.
\newblock \emph{arXiv:2305.19118}.

\bibitem[Liang et al.(2024)]{liang2024cotpot}
Liang et al. 2024.
\newblock {CoT}-{PoT} disagreement as an error detection signal.
\newblock \emph{arXiv preprint}.

\bibitem[Liu et al.(2021)]{liu2021logiqa}
Jian Liu, Leyang Cui, Hanmeng Liu, Dandan Huang, Yile Wang, and Yue Zhang. 2021.
\newblock {LogiQA}: A challenge dataset for machine reading comprehension with logical reasoning.
\newblock In \emph{Proc. IJCAI}.

\bibitem[Lyu et al.(2023)]{lyu2023faithful}
Qing Lyu, Shreya Havaldar, Adam Stein, Li Zhang, Delip Rao, Eric Wong, Marianna Apidianaki, and Chris Callison-Burch. 2023.
\newblock Faithful chain-of-thought reasoning.
\newblock In \emph{Proc. IJCNLP-AACL}.

\bibitem[Madaan et al.(2023)]{madaan2023selfrefine}
Aman Madaan, Niket Tandon, Prakhar Gupta, Skyler Hallinan, Luyu Gao, Sarah Wiegreffe, Uri Alon, Nouha Dziri, Shrimai Prabhumoye, Yiming Yang, Shashank Gupta, Bodhisattwa~Prasad Majumder, Katherine Hermann, Sean Welleck, Amir Yazdanbakhsh, and Peter Clark. 2023.
\newblock Self-refine: Iterative refinement with self-feedback.
\newblock In \emph{Proc. NeurIPS}.

\bibitem[Ni et al.(2023)]{ni2023lever}
Ansong Ni, Srini Iyer, Dragomir Radev, Veselin Stoyanov, Wen-tau Yih, Sida~I. Wang, and Xi~Victoria Lin. 2023.
\newblock {LEVER}: Learning to verify language-to-code generation with execution.
\newblock In \emph{Proc. ICML}.

\bibitem[Olausson et al.(2023)]{olausson2023linc}
Theo~X. Olausson, Alex Gu, Ben Lipkin, Cedegao~E. Zhang, Armando Solar-Lezama, Joshua~B. Tenenbaum, and Roger Levy. 2023.
\newblock {LINC}: A neurosymbolic approach for logical reasoning by combining language models with first-order logic provers.
\newblock In \emph{Proc. EMNLP}.

\bibitem[Pan et al.(2023)]{pan2023logiclm}
Liangming Pan, Alon Albalak, Xinyi Wang, and William~Yang Wang. 2023.
\newblock {Logic-LM}: Empowering large language models with symbolic solvers for faithful logical reasoning.
\newblock In \emph{Findings of EMNLP}.

\bibitem[PROVE(2024)]{prove2024}
PROVE. 2024.
\newblock Program verification for error detection in reasoning.
\newblock \emph{arXiv preprint}.

\bibitem[Rein et al.(2024)]{rein2024gpqa}
David Rein, Betty~Li Hou, Asa~Cooper Stickland, Jackson Petty, Richard~Yuanzhe Pang, Julien Dirani, Julian Michael, and Samuel~R. Bowman. 2024.
\newblock {GPQA}: A graduate-level Google-proof Q\&A benchmark.
\newblock In \emph{Proc. NAACL}.

\bibitem[Sun et al.(2024)]{sun2024adaptive}
Sun et al. 2024.
\newblock Adaptive representation selection for reasoning.
\newblock \emph{arXiv preprint}.

\bibitem[Wang et al.(2023)]{wang2023selfconsistency}
Xuezhi Wang, Jason Wei, Dale Schuurmans, Quoc~V. Le, Ed~H. Chi, Sharan Narang, Aakanksha Chowdhery, and Denny Zhou. 2023.
\newblock Self-consistency improves chain of thought reasoning in language models.
\newblock In \emph{Proc. ICLR}.

\bibitem[Weng et al.(2023)]{weng2023vericot}
Weng et al. 2023.
\newblock {VeriCoT}: Verification-based chain-of-thought.
\newblock \emph{arXiv preprint}.

\bibitem[Ye et al.(2024)]{ye2024satlm}
Xi Ye, Qiaochu Chen, Isil Dillig, and Greg Durrett. 2024.
\newblock {SatLM}: Satisfiability-aided language models using declarative prompting.
\newblock In \emph{Proc. NeurIPS}.

\bibitem[Zhong et al.(2025)]{zhong2025roundtrip}
Zhong et al. 2025.
\newblock Roundtrip verification for diagnosis-guided repair.
\newblock \emph{arXiv preprint}.

\end{thebibliography}

\end{document}
